\subsection{Primes}
\textbf{Wejście:} Liczba \( p \) \\
\textbf{Problem:} Czy \( p \) jest liczbą pierwszą? \\
\textbf{Klasa złożoności:} P

Primes \( \in \) coNP: \\
Jeśli otrzymamy od wyroczni \( d \) potencjalny dzielnik \( p \), możemy w wielomianowym czasie sprawdzić, czy \( d \) rzeczywiście jest dzielnikiem i odpowiedzieć NIE na pytanie o pierwszość.

Primes \( \in \) NP: \\
Korzystamy z tego, że \( \integer_p^* \) ma rząd równy \( p-1 \) wtedy i tylko wtedy, gdy \( p \) jest pierwsze. \\
Jeśli otrzymamy od wyroczni \( g \), generator \( \integer_p^* \) i rozkład \( p-1 = p_1^{\alpha_1} \cdots p_s^{\alpha_s} \), możemy:
\begin{itemize}
	\item rekurencyjnie sprawdzić pierwszość \( p_i \) i poprawność rozkładu,
	\item sprawdzić, czy \( g^{p-1} = 1 \pmod{p} \),
	\item sprawdzić, czy \( g^{\frac{p-1}{p_i}} \neq 1 \pmod{p} \).
\end{itemize}
Jeżeli wszystkie z powyższych warunków są spełnione, to odpowiedzią jest TAK.

Primes \( \in \) BPP: \\
Dowodem jest algorytm Millera-Rabina, który myli się przy odpowiedzi TAK z prawdopodobieństwem nie większym niż \( \frac{1}{2} \).

Primes \( \in \) P: \\
Dowodem jest algorytm Agrawala-Kayala-Saxena (AKS).

\subsection{Factoring}
\textbf{Wejście:} Liczby \( n,\; k \) \\
\textbf{Problem:} Czy istnieje nietrywialny dzielnik \( n \) mniejszy lub równy \( k \)? \\
\textbf{Klasa złożoności:} NP \( \cap \) coNP?

Factoring \( \in \) NP: \\
Jeśli otrzymamy od wyroczni \( 2 \leq d \leq k \) potencjalny dzielnik \( n \), możemy w wielomianowym czasie sprawdzić, czy \( d \) rzeczywiście jest dzielnikiem i odpowiedzieć TAK.

Factoring \( \in \) coNP: \\
Otrzymujemy od wyroczni rozkład liczby \( n \) na czynniki pierwsze \( n = p_1^{\alpha_1} \cdots p_s^{\alpha_s} \). Po sprawdzeniu poprawności rozkładu odpowiadamy \( (p_1 \leq k) \) ? TAK : NIE, gdzie \( p_1 \) to najmniejszy z~dzielników.

Na ten moment nie wiadomo nic więcej.

\subsection{Algorytm RSA}
Algorytm RSA jest asymetryczny, czyli wykorzystuje jawną funkcję szyfrującą \( E \) (klucz publiczny) i tajną funkcję deszyfrującą \( D \) (klucz prywatny).
Najważniejszym założeniem jest, że nie potrafimy efektywnie odtworzyć klucza prywatnego z publicznego. Dodatkowo zakładamy, że problem Factoring nie jest w klasie BPP.

\begin{greyframe}
	Algorytm RSA:
	\begin{enumerate}
		\item Oblicz \( N = p \cdot q \) dla wybranych liczb pierwszych \( p \), \( q \). \\
		      Wtedy \( \varphi(N) = (p-1) \cdot (q-1) \).
		\item Wybierz \( e \) względnie pierwsze z \( \varphi(N) \). Oblicz \( d \) takie, że  \( e\cdot d = 1 \pmod{\varphi(N)} \).
		\item Ustal klucz publiczny \( (e, N) \).
		\item Zaszyfruj wiadomość \( x \): \( E(x) = x^e \pmod{N} \)
		\item Odszyfruj wiadomość: \( D(x) = x^d \pmod{N} \)
	\end{enumerate}
\end{greyframe}

Ad 1. Zamiast \( \varphi(N) \) można użyć (rozwiązanie często stosowane) funkcji Carmichaela
\[
	\lambda(N) = \lcm(p - 1, q - 1)
\]
Funkcja \( \lambda \) znajduje najmniejszą liczbę całkowitą \( x \) taką, że \( x^{\lambda(N)} = 1 \pmod{N} \) dla każdego \( x \) względnie pierwszego z \( N \).
Wybór \( \lambda(N) \) zamiast \( \varphi(N) \) nie zmniejsza bezpieczeństwa, a daje zazwyczaj mniejszy wykładnik. \\
Ad 2. Do znalezienia liczby \( d \) można użyć rozszerzonego algorytmu Euklidesa:
\[
	e \cdot s + \varphi(N) \cdot t = 1 \implies e \cdot s = 1 \pmod{\varphi(N)}
\]
Dlaczego deszyfrowanie działa? \\
Rząd grupy multiplikatywnej modulo \( N \) jest równy \( \varphi(N) \), więc dla każdego \( x \) względnie pierwszego z \( N \) zachodzi
\[
	(x^{e})^d = x^{ed} = x^{k \cdot \varphi(N) + 1} = x,
\]
ponieważ \( e \cdot d = 1 \pmod{\varphi(N)} \), a na mocy twierdzenia Lagrange'a \( x^{\varphi(N)} = 1 \).

Szyfrowanie RSA można złamać, jeśli wyznaczy się liczby \( p \) i \( q \), czyli pozna się rozkład \( N \) na czynniki pierwsze.
Można wtedy obliczyć \( \varphi(N) \) i \( d \) algorytmem Euklidesa. Dlatego tak ważne jest założenie o dużej trudności problemu Factoring.
W tym dniu (20.06.2026), nie ma ani dowodu na to, że wielomianowy algorytm rozkładu na czynniki pierwsze nie istnieje, ani dowodu NP-zupełności problemu Factoring.

\subsection{Efektywna metoda znalezienia dużej liczby pierwszej}
Dla zadanej liczby \( k \) chcemy znaleźć liczbę pierwszą \( p \in \brackets{2^k, 2^{k+1}-1} \).

Liczby pierwsze są dość gęsto upakowane -- w przedziale od 1 do \( n \) jest około \( \frac{n}{\log n} \) liczb pierwszych.
Zatem w przedziale \( \brackets{2^k, 2^{k+1}-1} \) jest ich około
\[
	\frac{2^{k+1} - 1}{k+1} - \frac{2^{k} - 1}{k} = \frac{2^k\pars{1-\frac{1}{k}} + \frac{1}{k}}{k+1} = \Theta\pars{\frac{2^k}{k}}
\]

Możemy więc po prostu losować liczbę całkowitą \( p \) z przedziału \( \brackets{2^k, 2^{k+1}-1} \), dopóki nie trafi się liczba pierwsza.
Skoro w przedziale jest \( \Theta\pars{2^k} \) liczb całkowitych, z których \( \Theta\pars{\frac{2^k}{k}} \) jest liczbami pierwszymi, to w oczekiwaniu po \( \bigO(k) \) losowaniach trafimy na liczbę pierwszą.
Sprawdzanie, czy liczba jest pierwsza, wykonuje się w czasie \( \bigO\pars{k^3} \) (Millerem-Rabinem), więc oczekiwana złożoność algorytmu losowania liczby pierwszej to \( \bigO\pars{k^4} \).